% ============================================================================
% RETIRED 2026-08-25 -- NOT PART OF THE SUBMISSION.
% The author elected to eliminate supplementary material entirely. Everything
% retained is now folded into drafts/v6_draft.tex:
%   S1 Table  -> manuscript Table 1 (verbatim item wording)
%   S7 Table  -> manuscript Table 2 (missing-data decomposition)
%   S6 prose  -> Limitations (both skip-logic errors, described in full)
%   S2/S3     -> Results, Facilitator practice (expanded prose, no table)
%   S5        -> Limitations (both alphas stated inline)
%   S9        -> Data availability (reproducibility statement)
% Deliberately dropped: S3 BH family tables (pointers removed, procedure and
% every q-value kept in the main text), S1 coding notes, the preamble
% conventions, S4 proportional-odds model, S8 qualitative scaffold (held for a
% separate manuscript). Kept here as the analysis record only -- do not submit,
% and do not cite from the manuscript.
% ============================================================================
% supplement.tex — Supplementary Material for "Ceremonial Entheogen Use in Puerto Rico"
% Phase 5, round 2 revision (accompanies drafts/v4_draft.tex). Holds everything moved out of the
% main text to meet the JPD 4,000-word cap.
% Numbers trace to outputs/report_v2.md except the facilitator-branch distributions, the
% epsilon-squared values added in Phase 5 round 1, and the permutation p-value and median
% correction added in Phase 5 round 2, all direct recomputations from
% data/results_7_15_2026.xlsx via analysis/data_prep.load() / analysis/stats_helpers.py;
% instrument wording traces to manuscript/sources/instrument_documentation.md (extracted
% from data/koboxls.xlsx).

\documentclass[12pt]{article}

\usepackage[utf8]{inputenc}
\usepackage[T1]{fontenc}
\usepackage{times}
\usepackage[margin=1in]{geometry}
\usepackage{graphicx}
\usepackage{booktabs}
\usepackage{longtable}
\usepackage{amsmath}
\usepackage{array}
\emergencystretch=3em

\renewcommand{\thetable}{S\arabic{table}}
\renewcommand{\thesection}{S\arabic{section}}

\title{Supplementary Material\\
\large Ceremonial Entheogen Use in Puerto Rico: A Descriptive Survey of Participants and a Small
Facilitator Subsample}
\date{}

\begin{document}
\maketitle

\noindent
All results below derive from the same analytic sample described in the main text ($N=73$ respondents who
identified a role; participant branch $n=72$ self-identified of whom 67 were routed into the participant section
(3 dual-role respondents were stopped by the continuation gate and 2 answered ``prefer not to answer''
to the facilitator role item, satisfying neither arm of the section's relevance condition), facilitator branch $n=8$, 7 dual-role). The structural-skip and ``prefer not to answer'' conventions
are those stated in the Methods and are applied uniformly here. Counts are exact throughout, including
cells below five: as the Methods explain, the disclosure protection is the anonymous, aggregate design
rather than cell masking. The bias-corrected $\varepsilon^2$ formula can return small negative values when
an effect is indistinguishable from zero; these are displayed truncated to 0.00 and marked $\dagger$,
applied uniformly.

\section{Instrument: verbatim wording and scale anchors}

The questionnaire was deployed as a single bilingual KoboToolbox XLSForm with Spanish as the default
language and English selectable through the platform's language switcher. Every item carries both a Spanish
and an English label in the same form; this is a native bilingual instrument rather than two separate
questionnaires or a post-hoc translation. English wording is reproduced below.

\begin{table}[htbp]
\centering
\caption{Verbatim wording and anchors for the items carrying the main text's principal results.}
\begin{tabular}{p{3cm}p{6cm}p{5cm}}
\toprule
Item (construct) & Wording & Response options \\
\midrule
Satisfaction \newline (primary outcome) &
``From 0 to 6, how pleasant was the ceremony?'' &
0 Very unpleasant; 1 Unpleasant; 2 Somewhat unpleasant; 3 Neither pleasant nor unpleasant; 4 Somewhat
pleasant; 5 Pleasant; 6 Very pleasant; Prefer not to answer \\
\addlinespace
Self-directed preparation \newline (primary exposure) &
``Aside from the facilitator's instructions or recommendations, do you personally engage in preparation for
these ceremonies?'' &
Yes; Sometimes; No; Prefer not to answer \\
\addlinespace
In-ceremony comfort and safety &
``During the ceremony, how would you describe your overall sense of comfort and safety?'' &
0 I felt very uncomfortable or unsafe; 1 At times I felt somewhat uncomfortable or unsafe; 2 I felt mostly
comfortable and safe; 3 I felt completely comfortable and safe; Prefer not to answer \\
\addlinespace
Trust in facilitator &
``How much did you trust the person who facilitated the ceremony?'' &
0 Not at all; 1 A little; 2 Moderately; 3 A lot; 4 Completely; Prefer not to answer \\
\addlinespace
Non-consensual physical contact \newline (lifetime) &
``Have you ever experienced physical touch during a ceremony that felt unexpected or that you did not
clearly agree to beforehand? This could include moments where you were unsure, uncomfortable, or did not
remember giving permission.'' &
Yes; No; I'm not sure; Prefer not to answer \\
\addlinespace
Non-consensual physical contact \newline (most recent ceremony) &
``During your most recent ceremony, did you experience any physical touch that felt unexpected or that you
didn't clearly agree to beforehand?'' &
Yes; No; I'm not sure; Prefer not to answer \\
\bottomrule
\end{tabular}
\end{table}

\noindent
Three notes on coding. First, the preparation item measures presence and frequency of the respondent's
\emph{own} preparation, deliberately separate from any instructions the facilitator gave; a companion
open-text field captured what that preparation consisted of. The construct is therefore heterogeneous by
design rather than by instrument failure. It also asks about ``these ceremonies'' as a general practice,
while the satisfaction outcome and the two facilitator-practice items ask about the most recent ceremony
alone --- a referent difference that the main text discloses wherever the two are related. Second, the raw
choices sheet lists the comfort-and-safety
options best-first, the reverse of the ascending 0--3 code used throughout the analysis. The analysis
re-orders these by semantic content so that a higher code always means a more positive response, consistent
with every other ordinal scale in the study. A reviewer cross-checking the raw XLSForm should expect this
difference.

\section{Facilitator self-reported practice}

Eight respondents identified as facilitators, and \textbf{seven of the eight also identified as
participants}. Every facilitator who answered a substantive facilitator item is one of those seven, so this
branch is largely a subset of the participant branch rather than a separate set of people, and the tables
below are \emph{not} independent of the participant-side results in the main text. What remains unlinked is
the \emph{ceremony}: no ceremony or facilitator identifier connects any participant's report to the
facilitator who ran that ceremony, and for a dual-role respondent the ceremonies led are different events
from the ones attended. No group comparison at this base is meaningful and none is presented; these tables
are descriptive context only and must not be read against any participant-reported item as though the two
described the same ceremonies.

\noindent
\textbf{A note on denominators.} Eight respondents entered this branch but only seven answered any
facilitator item; the eighth (the sole facilitator-only respondent) answered none. Every percentage in both
tables below is therefore computed on the seven who answered, not on the base of eight.

\begin{table}[htbp]
\centering
\caption{Facilitator training, screening approach, and distress response (select-multiple items; base
$n=7$ answering; percentages of that base need not sum to 100\%).}
\begin{tabular}{p{2.8cm}p{7.2cm}rr}
\toprule
Item & Category & $n$ & \% answering \\
\midrule
Training / experience & Personal experience and self-study & 6 & 86 \\
                      & Mentorship or apprenticeship learning & 4 & 57 \\
                      & Family or ancestral tradition & 3 & 43 \\
                      & Formal certifications & 0 & 0 \\
                      & Academic courses or workshops & 0 & 0 \\
\addlinespace
Screening approach    & Initial conversation or interview with the person & 6 & 86 \\
                      & Asks about relevant medical history & 4 & 57 \\
                      & Asks about current medication & 4 & 57 \\
                      & Asks about mental health challenges & 3 & 43 \\
                      & Welcomes anyone who feels called to participate & 2 & 29 \\
                      & Other & 2 & 29 \\
\addlinespace
Response to a         & Singing, playing music, or using sound & 5 & 71 \\
difficult experience  & Grounding exercises (breathing, sensory focus) & 4 & 57 \\
                      & Calming or supportive words & 3 & 43 \\
                      & Physical touch (with consent or as culturally appropriate) & 2 & 29 \\
                      & Other & 2 & 29 \\
                      & ``Trip killer'' medication & 1 & 14 \\
                      & Has never been in this situation & 1 & 14 \\
\bottomrule
\end{tabular}
\end{table}

\begin{table}[htbp]
\centering
\caption{Facilitator practice context, emergency planning, witnessed crisis signs, and readiness composite.
Base is the same $n=7$ answering used in the preceding table.}
\begin{tabular}{p{4.6cm}p{5.4cm}rr}
\toprule
Item & Category & $n$ & \% answering \\
\midrule
Gives preparation instructions & Yes & 6 & 86 \\
                             & Prefer not to answer & 1 & 14 \\
\addlinespace
Typical environment          & Outdoor space or near nature & 6 & 86 \\
                             & In my home & 1 & 14 \\
\addlinespace
Dosing approach              & Participants decide their own dose & 3 & 43 \\
                             & Other & 2 & 29 \\
                             & Conservative initial dose, then more if needed & 1 & 14 \\
                             & A single dose for everyone & 1 & 14 \\
\addlinespace
Emergency plan               & Informal or unwritten approach & 4 & 57 \\
                             & Depends on the situation & 3 & 43 \\
                             & Formal / written plan & 0 & 0 \\
\addlinespace
Crisis signs witnessed       & They appear disoriented & 3 & 43 \\
(select-multiple)            & They look overwhelmed & 3 & 43 \\
                             & Fainting or dizziness & 3 & 43 \\
                             & Difficulty breathing & 3 & 43 \\
                             & Panic attacks & 3 & 43 \\
                             & I'm not sure / it depends & 3 & 43 \\
                             & No response or loss of consciousness & 2 & 29 \\
                             & Erratic or unusual behavior & 2 & 29 \\
                             & Rapid heartbeat or physical agitation & 2 & 29 \\
                             & Signs of self-harm or suicidal thoughts & 2 & 29 \\
                             & Aggressive or threatening behavior & 2 & 29 \\
                             & Allergic reactions & 2 & 29 \\
                             & Paranoia or extreme fear & 2 & 29 \\
                             & Other & 2 & 29 \\
                             & Reliving or re-experiencing past trauma & 1 & 14 \\
                             & Excessive purging (vomiting, diarrhea) & 0 & 0 \\
\addlinespace
Readiness composite          & Score 3.5 of 4 & 6 & 86 \\
                             & Score 2.5 of 4 & 1 & 14 \\
\bottomrule
\end{tabular}
\end{table}

\noindent
The facilitator branch is close to constant: every facilitator answering reported some form of training and
some screening approach, and none selected the formal emergency-plan option. With that little variance at
$n=7$, neither a reliability statistic nor an association test is informative. Three of the seven reported
several crisis signs each, two reported a single sign, one answered only ``I'm not sure'' and one only
``other'' --- so the counts above concentrate in three respondents and must not be read as rates.
Association tests between witnessed crisis signs and facilitator-reported protocol items---the one place in
the study where both variables sit on the same side of the survey and no linkage gap exists---returned
undefined odds ratios: Fisher's exact $p=0.143$ against the distress-response protocol and $p=0.429$
against the emergency plan, with two further tests unrunnable for insufficient non-degenerate data. Every
$2\times2$ table here has at least one margin of 2 or fewer. These results establish an analytic approach
for a future, larger facilitator sample; they are not evidence for or against an association.

\section{Multiplicity correction: all four test families}

The preparation-to-satisfaction comparison, designated in the analysis plan as the primary comparison
before any test was run, is reported uncorrected in the main text and is excluded from every family below.
The study was not preregistered. All other bivariate tests were assigned in advance to one of four
thematically organized families and corrected jointly within family by the Benjamini--Hochberg
false-discovery-rate procedure. Two variables, non-consensual contact and screening received, recur as
predictor or outcome across three of the four families (see the safety-domain and legal-knowledge families
below), so family-wise correction does not bound the error rate for conclusions about either variable across
every test in which it appears; no q-value in any family is below 0.05, so this does not change any result.
\textbf{No test in any family was significant at $q<0.05$.}

\begin{table}[htbp]
\centering
\caption{Benjamini--Hochberg correction, bivariate family. $\dagger$ marks an $\varepsilon^2$ whose
bias-corrected value was negative and is displayed truncated to 0.00. $\ddagger$: a tie-aware permutation
check (10,000 label permutations, seed 20260715) gave $p=0.010$ for this test, the only one in this family
reaching uncorrected significance; its ``no''-background group contributes only 3 non-missing values.}
\begin{tabular}{lrrr}
\toprule
Test & raw $p$ & BH $q$ & $\varepsilon^2$ \\
\midrule
Preparation (3-level) $\rightarrow$ satisfaction (KW) & 0.120 & 0.300 & 0.04 \\
Comfort/safety by prior knowledge of facilitator background (KW) & 0.326 & 0.407 & 0.00 \\
Comfort/safety by screening received (KW) & 0.301 & 0.407 & 0.01 \\
Trust in facilitator by prior knowledge of background (KW) & $0.017^{\ddagger}$ & 0.085 & 0.10 \\
Trust in facilitator by screening received (KW) & 0.679 & 0.679 & $0.00^{\dagger}$ \\
\bottomrule
\end{tabular}
\end{table}

\begin{table}[htbp]
\centering
\caption{Benjamini--Hochberg correction, substance-stratified family.}
\begin{tabular}{lrrr}
\toprule
Test & raw $p$ & BH $q$ & effect size \\
\midrule
Satisfaction by substance group (KW) & 0.083 & 0.166 & $\varepsilon^2=0.05$ \\
Non-consensual contact by substance group (permutation $\chi^2$) & 0.486 & 0.486 & $V=0.17$ \\
\bottomrule
\end{tabular}
\end{table}

\begin{table}[htbp]
\centering
\caption{Benjamini--Hochberg correction, safety-domain family. The non-consensual contact outcome is the
lifetime item, while the three participant-side exposures crossed against it all ask about the most recent
ceremony; 3 of the 5 respondents reporting the event ever placed it at an earlier ceremony, so for those
three the exposure and the outcome describe different events. This referent mismatch applies to every row
in the upper block and is the main reason these tests are hypothesis-generating only.}
\begin{tabular}{p{8.2cm}rrl}
\toprule
Test & raw $p$ & BH $q$ & effect size \\
\midrule
Non-consensual contact $\times$ prior knowledge of facilitator background & 0.047 & 0.188 & $V=0.38$ \\
Non-consensual contact $\times$ screening received & 0.094 & 0.188 & $V=0.32$ \\
\addlinespace
Witnessed crisis signs $\times$ distress-response protocol (facilitator branch) & 0.143 & 0.191 & OR undefined \\
Witnessed crisis signs $\times$ emergency plan (facilitator branch) & 0.429 & 0.429 & OR undefined \\
\bottomrule
\end{tabular}
\end{table}

\begin{table}[htbp]
\centering
\caption{Benjamini--Hochberg correction, legal-knowledge family. Predictors are three participant-side
legal items; outcomes are disclosure to a health professional, screening received, and non-consensual
contact. All tests Kruskal--Wallis. $\dagger$ marks a truncated $\varepsilon^2$, as above.}
\begin{tabular}{lrrr}
\toprule
Test & raw $p$ & BH $q$ & $\varepsilon^2$ \\
\midrule
General legal knowledge by disclosure to a health professional & 0.258 & 0.581 & 0.01 \\
Most-recent-ceremony legal knowledge by disclosure & 0.164 & 0.495 & 0.03 \\
Most-recent-ceremony legal worry by disclosure & 0.031 & 0.283 & 0.09 \\
General legal knowledge by screening received & 0.412 & 0.741 & $0.00^{\dagger}$ \\
Most-recent-ceremony legal knowledge by screening received & 0.659 & 0.741 & $0.00^{\dagger}$ \\
Most-recent-ceremony legal worry by screening received & 0.914 & 0.914 & $0.00^{\dagger}$ \\
General legal knowledge by non-consensual contact & 0.617 & 0.741 & $0.00^{\dagger}$ \\
Most-recent-ceremony legal knowledge by non-consensual contact & 0.165 & 0.495 & 0.03 \\
Most-recent-ceremony legal worry by non-consensual contact & 0.611 & 0.741 & $0.00^{\dagger}$ \\
\bottomrule
\end{tabular}
\end{table}

\noindent
The only legal-domain comparison reaching uncorrected significance was most-recent-ceremony legal worry by
disclosure to a health professional ($H=6.92$, $p=0.031$, $\varepsilon^2=0.09$): median 1.5 (IQR 0--2) among
the 8 reporting disclosure against median 0 (IQR 0--0) among the 47 reporting none. It did not survive
correction ($q=0.283$). Two further legal items were not testable: a general legal-worry item reached only
six respondents because of the skip-logic error described in \S\ref{sec:logic}, and a facilitator
legal-knowledge item sits on the unlinked facilitator branch with no participant-side counterpart.

\section{Exploratory proportional-odds model}

This model adjusts preparation for age and facilitator trust simultaneously. The full covariate set
specified in the original protocol was not estimable: gender cells were too thin for a four-level
categorical predictor, and screening overlapped substantively with the trust and background measures
already in the model. A prespecified three-predictor fallback was used instead. The outcome was collapsed
to three ordered levels for this model only, because the raw seven-level scale has cells as small as one to
three respondents at the low end and risks quasi-complete separation at this sample size; every other table
in this paper retains the full 0--6 scale.

\begin{table}[htbp]
\centering
\caption{Proportional-odds model for reported satisfaction. Complete-case $n=62$ of the 63 participants
with a non-missing outcome. Collapsed outcome: low/mixed (0--3) $n=9$; mostly positive (4--5) $n=19$; very
positive (6) $n=34$.}
\begin{tabular}{lrrr}
\toprule
Predictor & OR & 95\% CI & $p$ \\
\midrule
Age (years) & 0.99 & 0.94--1.03 & 0.527 \\
Any preparation vs.\ none & 5.21 & 0.84--32.41 & 0.077 \\
Trust in facilitator (0--4) & 2.20 & 1.06--4.53 & 0.033 \\
\bottomrule
\end{tabular}
\end{table}

\noindent
Omnibus likelihood-ratio test against an intercept-and-cutpoint-only model: $\chi^2(3)=8.68$, $p=0.034$;
McFadden pseudo-$R^2=0.07$. \textbf{Power caveat.} Sixty-two complete cases support five estimated
parameters (three slopes and two cutpoints), roughly 12 cases per parameter, below conventional rules of
thumb for stable ordinal-regression estimates. The proportional-odds assumption was not formally tested,
the sample being too small to test it reliably, and is simply assumed. Every estimate in this table is
exploratory and hypothesis-generating, and should not be presented without this caveat.

\section{Psychometric evidence for the instrument}

The instrument was purpose-built, so this is its first reliability evidence. The original protocol's
language implies multi-item Likert batteries, but the form contains none: every Likert-shaped option set
maps to at most two items, and most two-item pairs split across the unlinked facilitator and participant
branches. Two genuine same-respondent pairs are testable.

\begin{table}[htbp]
\centering
\caption{Reliability of the two testable same-respondent item pairs.}
\begin{tabular}{lrrr}
\toprule
Scale & Cronbach's $\alpha$ & 95\% CI & $n$ \\
\midrule
Trust in facilitator + in-ceremony comfort/safety (2 items) & 0.58 & 0.31--0.74 & 64 \\
Legal knowledge, general + most-recent ceremony (2 items) & 0.86 & 0.75--0.92 & 55 \\
\bottomrule
\end{tabular}
\end{table}

\noindent
With only two items, $\alpha$ is a monotonic transform of the pairwise correlation and the item-rest
correlations are not independent numbers. The corresponding Spearman correlations are $\rho=0.39$
($p=0.002$, $n=64$) for the trust-and-comfort pair and $\rho=0.79$ ($p<0.001$, $n=55$) for the
legal-knowledge pair; the latter two items required harmonization onto a shared 0--4 code because the
most-recent-ceremony version omits the bottom category its general counterpart carries.

A third candidate pair was untestable: a perceived-life-impact item had zero variance, with all 67
respondents answering selecting the same option, so it cannot correlate with anything and $\alpha$ is
undefined. This is reported as a data-quality finding, not a scale result.

\section{Instrument skip-logic errors}
\label{sec:logic}

Two errors were identified in the fielded form and are reported so that any future deployment corrects
them.

First, a follow-up item asking why a respondent felt uncomfortable or unsafe received zero responses. Its
relevance condition required the comfort-and-safety variable to equal two different values simultaneously,
making the item structurally unreachable rather than merely unpopular. No respondent could have answered
it.

Second, a general legal-worry item was gated on the general legal-knowledge item equaling ``prefer not to
answer''---the opposite of what a parallel item would require---so only six respondents ever reached it,
while its most-recent-ceremony counterpart was asked broadly ($n=57$). The two are therefore neither
comparable nor parallel, and the general item is excluded from all analyses.

\section{Missing-data distribution, attention checks, and duplicate adjudication}

Table~S10 decomposes missingness item by item over the analytic sample ($N=73$) into four mutually
exclusive states: answered substantively, answered with a retained non-response category (``prefer not
to answer'' / ``not sure''), left blank while routed, and never shown by the skip logic. Routing is
evaluated from the group-level \emph{relevant} expressions transcribed from the XLSForm, not inferred
from which cells happen to be empty --- the distinction matters, because the two are not the same set.
The four states partition the analytic sample by construction, and the pipeline asserts that they do.

The decomposition shows that missingness on the participant items is overwhelmingly structural. Of the
72 participant-side respondents, 5 were never routed; among the 67 who were, item nonresponse peaks at 3
respondents (4.5\%) on the satisfaction outcome and is zero on six of the fifteen participant items.
Retained non-response is more common than blankness and is concentrated on one item: 12 of the 67
answered ``not sure'' or ``prefer not to answer'' on whether the facilitator asked screening questions,
which is itself a substantive finding about how visible screening is to participants rather than a
missing-data problem. Nothing here reaches a level at which imputation would be warranted, and none was
performed.

Attention checks were branch-specific rather than uniform: facilitators were judged on three items and
participants on four. Two respondents answered at least one attention item incorrectly. Per protocol these
rows were flagged and retained, not dropped.

The survey platform never populated its duplicate-detection field and recorded no IP or device information,
so duplicate submissions were screened manually on the respondent-chosen nickname and password. Three
signals were examined: an exact nickname-and-password match, a repeated nickname alone, and a repeated
password alone. Only the first indicates a duplicate submission rather than coincidental reuse of a common
word.

Exactly one exact-match pair exists across the 100 submissions. Its earlier member is an abandoned blank
submission two days before the completed one---the same respondent trying twice, abandoning, and
completing on the second attempt---and was already excluded by the pre-existing abandonment rule.
\textbf{Zero analytic rows require exclusion on duplicate grounds}, and no sensitivity re-run comparing
adjudicated against unadjudicated samples is possible or necessary, because the two samples are identical.

Nickname-only repeats (9 rows) and password-only repeats (12 rows) were inspected individually. The
repeated strings are common Spanish words that respondents in this thematic context would plausibly choose
independently---\emph{Amor}, \emph{S\'i}, \emph{Secreta}, \emph{Hongos}, \emph{Luna}---and in every case
where a repeat might conceivably indicate the same person, all members already fall outside the analytic
sample as abandoned or screened-out submissions. No pair has both members inside the analytic sample with a
same-person signal strong enough to justify exclusion.

\begin{table}[htbp]
\centering
\small
\caption{Missing-data decomposition by item, analytic sample ($N=73$). The four states are mutually
exclusive and sum to 73 on every row. \emph{Substantive} = a real response category; \emph{PNA/not sure} =
a retained non-response category (kept as its own category throughout, never dropped); \emph{Blank} =
routed to the item and left it empty, the only genuine item nonresponse; \emph{Not routed} = the skip
logic never showed the item. Routing conditions are transcribed from the XLSForm \emph{relevant}
expressions, not inferred from empty cells.}
\label{tab:missingness}
\begin{tabular}{lrrrr}
\toprule
Item & Substantive & PNA/not sure & Blank & Not routed \\
\midrule
\multicolumn{5}{l}{\emph{Sociodemographic branch --- routed $n=73$ of 73}} \\
\midrule
Age & 72 & 0 & 1 & 0 \\
Gender & 71 & 0 & 2 & 0 \\
Education & 72 & 0 & 1 & 0 \\
Economic situation & 70 & 2 & 1 & 0 \\
Country of residence & 72 & 0 & 1 & 0 \\
\midrule
\multicolumn{5}{l}{\emph{Participant branch --- routed $n=67$ of 73}} \\
\midrule
Substances, all ceremonies attended & 67 & 0 & 0 & 6 \\
Substance, most recent ceremony & 65 & 1 & 1 & 6 \\
Time since last ceremony & 65 & 1 & 1 & 6 \\
Place, most recent ceremony & 66 & 1 & 0 & 6 \\
Format, most recent ceremony & 62 & 3 & 2 & 6 \\
Duration, most recent ceremony & 65 & 0 & 2 & 6 \\
Self-directed preparation & 67 & 0 & 0 & 6 \\
Knew facilitator's background & 66 & 1 & 0 & 6 \\
Facilitator asked screening questions & 54 & 12 & 1 & 6 \\
Spoke to a health professional beforehand & 62 & 3 & 2 & 6 \\
Satisfaction, most recent ceremony & 63 & 1 & 3 & 6 \\
Trust in facilitator & 65 & 2 & 0 & 6 \\
In-ceremony comfort and safety & 65 & 1 & 1 & 6 \\
Non-consensual contact, ever & 66 & 1 & 0 & 6 \\
Non-consensual contact, most recent & 64 & 1 & 2 & 6 \\
\midrule
\multicolumn{5}{l}{\emph{Facilitator branch --- routed $n=8$ of 73}} \\
\midrule
Facilitator training background & 7 & 0 & 1 & 65 \\
Emergency plan & 7 & 0 & 1 & 65 \\
Crisis signs witnessed & 6 & 1 & 1 & 65 \\
\bottomrule
\end{tabular}
\end{table}

\section{Qualitative corpus and coding scaffold}

Thirty-seven open-text fields yielded 285 non-empty answers, mixed Spanish and English (220 answers, 77\%,
classified as Spanish; 54, 19\%, English; 11, 4\%, mixed or other, by a heuristic that characterizes the
corpus rather than validly classifying each item). The most-answered fields were an experience highlight
(58 answers), an opinion-on-safety item (54), a description of preparation guidance received (45), a
description of screening questions asked (38), and a description of self-directed preparation (30).

A seed keyword scan against the research team's prespecified categories returned legality (19 answers),
medical presence (19), facilitator vetting (14), education (13), and consent (1). That the consent category
is almost empty in a corpus where legality and facilitator vetting are the dominant themes is itself worth
noting for the successor instrument, given that 5 respondents reported non-consensual physical contact on
the closed-ended item.
The corpus is provided to the research team in long format with an accompanying manual-coding template.
This is a scaffold to support thematic analysis, not a substitute for it, and no interpretation of the
qualitative material is offered in the main text.

\section{Reproducibility}

Every number in the main text and in this supplement is generated by an executable analysis pipeline
(Python) that reads the raw survey export and the XLSForm directly; all response decodes derive from the
form rather than from hand-entered value maps. The pipeline regenerates the full statistical report in one
command. Code is available as described in the main text's data-availability statement.

\end{document}
